% paper/sections/08_pressure.tex
\section{圧力 Poisson 方程式（PPE）の離散化と求解}
\label{sec:pressure}

本章では，変密度 Projection 法（§~\ref{sec:projection_derivation}）から得られる
圧力 Poisson 方程式を段階的に構築し，それを高精度・効率的に解く枠組みを整える．
章の構成は以下の通りである：

\begin{enumerate}
  \item \textbf{§~\ref{sec:projection_derivation}（PPE の導出）}：
        変密度 IPC 法の Predictor--PPE--Corrector 3ステップを導出し，
        可変係数演算子 $\bnabla\cdot(1/\rho\,\bnabla p)$ が生じる理由を示す．
  \item \textbf{§~\ref{sec:ccd_poisson_matrix}（CCD 離散化）}：
        PPE を CCD 法で $\Ord{h^6}$ 精度に離散化し，
        DCCD によるチェッカーボード抑制と界面 Balanced-Force 整合を整える．
  \item \textbf{§~\ref{sec:defect_correction_main}（欠陥補正法）}：
        「高精度で評価し，低精度で解く」原理により $\Ord{h^6}$ 精度を保ったまま
        $\mathcal{O}(N^2)$ コストで求解する．
  \item \textbf{§~\ref{sec:field_extension}（Hermite 場延長法）}：
        Young--Laplace 不連続による CCD の Gibbs 振動を，
        Hermite 場延長法（HFE, $\Ord{h^6}$）で PPE 求解前に除去する．
  \item \textbf{§~\ref{sec:ppe_bc_stability}（境界条件と安定性）}：
        全周 Neumann 境界の零モード対策と高密度比における条件数を整理する．
\end{enumerate}

% ═══════════════════════════════════════════════════════════════════════════════
\subsection{変密度 Projection 法の段階的導出}
\label{sec:projection_derivation}
% ═══════════════════════════════════════════════════════════════════════════════

第~\ref{sec:collocate}章（§~\ref{sec:helmholtz}）では等密度場を仮定した Helmholtz 分解で
Projection 法の幾何学的意味を説明した．
本節ではこれを\textbf{変密度場}へ拡張し，
なぜ可変係数演算子 $\bnabla\cdot(1/\rho\,\bnabla p)$ が現れるかを3ステップで導出する．

本稿では van Kan（1986）の \textbf{IPC（Incremental Pressure Correction）法}を採用し，
予測ステップに前時刻圧力 $-\bnabla p^n$ を組み込む
（§~\ref{sec:time_int}，式~\eqref{eq:predictor_ab2_ipc}）．
これにより標準 Chorin 法の $\Ord{\Delta t}$ スプリッティング誤差が $\Ord{\Delta t^2}$ に低減される．

\subsubsection{Step 1：予測速度（Predictor）}

時刻 $t^{n+1}$ の離散 NS 方程式（無次元形）は：
\begin{equation}
  \rho^{n+1}\frac{\bu^{n+1} - \bu^n}{\Delta t}
  = \underbrace{-\rho^{n+1}(\bu^n\cdot\bnabla)\bu^n
    + \frac{1}{Re}\bnabla\cdot[\mu^{n+1}(\bnabla\bu^{n+1}+\bnabla^T\bu^{n+1})]
    + \bm{F}_\text{ext}^{n+1}}_{\displaystyle \equiv \bm{R}^{n+1}}
  - \bnabla p^{n+1}
  \label{eq:NS_full}
\end{equation}
IPC 法では前時刻圧力 $-\bnabla p^n$ を陽的に加えて\textbf{予測速度} $\bu^*$ を定義する：
\begin{equation}
  \rho^{n+1}\frac{\bu^* - \bu^n}{\Delta t} = \bm{R}^{n+1} - \bnabla p^n
  \quad\Longrightarrow\quad
  \bu^* = \bu^n + \frac{\Delta t}{\rho^{n+1}}\bigl(\bm{R}^{n+1} - \bnabla p^n\bigr)
  \label{eq:predictor}
\end{equation}
$\bu^*$ は一般に発散フリーではない（$\bnabla\cdot\bu^* \neq 0$）．

\subsubsection{Step 2：圧力補正（Corrector）}

式~\eqref{eq:NS_full} から式~\eqref{eq:predictor} を引くと，
圧力増分 $\delta p \equiv p^{n+1} - p^n$ と速度補正の関係が得られる：
\begin{equation}
  \rho^{n+1}\frac{\bu^{n+1} - \bu^*}{\Delta t} = -\bnabla(\delta p)
  \quad\Longrightarrow\quad
  \bu^{n+1} = \bu^* - \frac{\Delta t}{\rho^{n+1}}\bnabla(\delta p)
  \label{eq:corrector}
\end{equation}

\subsubsection{Step 3：発散フリー条件から PPE を導く}

$\bu^{n+1}$ に非圧縮条件 $\bnabla\cdot\bu^{n+1}=0$ を課す．
式~\eqref{eq:corrector} の両辺に発散を作用させると：
\begin{equation}
  \bnabla\cdot\!\left(\frac{1}{\rho^{n+1}}\bnabla(\delta p)\right)
  = \frac{1}{\Delta t}\,\bnabla\cdot\bu^*
  \label{eq:PPE_varrho}
\end{equation}
これが変密度 PPE（IPC 増分形式）である．
一様密度（$\rho=\mathrm{const}$）では $\bnabla^2(\delta p) = \rho\,\bnabla\cdot\bu^*/\Delta t$ に帰着する．

\begin{tcolorbox}[mybox, title={実装上の注意}]
\begin{itemize}
  \item \textbf{外力 $\bm{F}_\text{ext}^{n+1}$} には重力項と CSF 体積力
        $(\kappa^{n+1}\bnabla\psi^{n+1})/We$ の両方が含まれる
        （式~\eqref{eq:NS_dimless_cls}；CSF は Predictor に保持）．
  \item \textbf{$\bnabla p^n$ の評価前} に HFE（§~\ref{sec:field_extension}，Step 5c）で
        $p^n_{\mathrm{ext}}$ を構築する．
  \item \textbf{$\bnabla(\delta p)$ の評価前} に
        HFE（§~\ref{sec:field_extension}，Step 7 前処理）で $\delta p_{\mathrm{ext}}$ を構築する．
  \item \textbf{$\rho^{n+1}$} は CLS 移流（§~\ref{sec:algorithm}，Step 1）で事前に更新済み．
  \item \textbf{時間離散化：} 本節の導出は Backward Euler 形式で統一したが，
        対流項は実装で AB2（式~\eqref{eq:predictor_ab2_ipc}），
        粘性項は Crank--Nicolson（$\Ord{\Delta t^2}$，§~\ref{sec:advection}）を採用する．
\end{itemize}
\end{tcolorbox}

\noindent\textbf{PPE の離散化と求解方針：}
以下の演算子形式は圧力増分 $\delta p$ および圧力 $p$ のいずれにも適用可能であるため，
一般性のため $p$ で表記する．
式~\eqref{eq:PPE_varrho} の可変係数演算子を product-rule で展開すると：
\begin{equation}
  \bnabla\cdot\!\left(\frac{1}{\rho}\bnabla p\right)
  = \frac{1}{\rho}\bnabla^2 p - \frac{\bnabla\rho}{\rho^2}\cdot\bnabla p
  \label{eq:varrho_product_rule}
\end{equation}
この形式を CCD 法（$\Ord{h^6}$）で離散化する（§~\ref{sec:ccd_poisson_matrix}）．
得られた線形系は欠陥補正法（§~\ref{sec:defect_correction_main}）で効率的に解く．
FVM との等価性・調和平均の根拠は §~\ref{sec:fvm_face_coeff} 参照．
実装詳細（仮想時間刻み・ゲージ固定）は付録~\ref{sec:ppe_pseudotime} 参照．
